
\documentclass[10pt,a4paper,twocolumn,twoside]{rho-class/rho}
\definecolor{rhocolor}{rgb}{0.0, 0.2, 0.4} % dark blue accent (the template's built-in blue option)
\usepackage[english]{babel}
\usepackage{algorithm}
\usepackage{algpseudocode}
\usepackage{multirow}
\usepackage{pgfplots}
\pgfplotsset{compat=1.18}
\setlength{\columnsep}{0.25in}
% Visible marker for content that depends on Atlas/Vanda runs; none may remain at submission.
\newcommand{\pending}[1]{\textcolor{red}{\textbf{[PENDING: #1]}}}
\setlength{\intextsep}{6pt}
\setlength{\textfloatsep}{6pt}
\usepackage{titlesec}
\titlespacing*{\paragraph}{0pt}{0.2ex}{0.5em}
\makeatletter
\titlespacing*{\subsubsection}{0pt}{1ex}{2pt}
\makeatother
%----------------------------------------------------------
% Title
%----------------------------------------------------------

\title{DSA5208 Project 3 Report : MPI+OpenMP Parallelization of Barnes--Hut N-Body Simulation}

\author{\textbf{Group Members: }Yan Yixin (A0327039A), Zhang Yue (A0333668U), Leon Yiliang Qian (A0355825W)}

%----------------------------------------------------------
% Abstract and Keywords
%----------------------------------------------------------

\begin{abstract}

\end{abstract}
\makeatletter
\renewcommand{\rhoabstract}{}
\makeatother

%----------------------------------------------------------
% Footer: page number centred at the bottom of EVERY page, as "n/N"
%----------------------------------------------------------
\makeatletter
% Running head: project title on every page that carries one
\fancyhead[RO,LE]{\footerfont DSA5208 Project 3}
\fancyhead[RE,LO]{}
\fancyfoot[C]{\footerfont\thepage/\pageref{LastPage}}
\fancypagestyle{firststyle}{
    \fancyhead[L]{}
    \fancyhead[C]{}
    \fancyhead[R]{}
    \fancyfoot[C]{\footerfont\thepage/\pageref{LastPage}}
    \fancyfoot[L]{
        {\ifx\@journalname\undefined\else\footerfont\bfseries\@journalname\hspace{8pt}\fi}
        {\ifx\@theday\undefined\else\footerfont\@theday\hspace{8pt}\fi}
        {\ifx\@vol\undefined\else\footerfont{Vol.\nobreakspace}\@vol\hspace{8pt}\fi}
        {\ifx\@no\undefined\else\footerfont{No.\nobreakspace}\@no\hspace{8pt}\fi}
        }
    \fancyfoot[R]{
        {\ifx\@journal\undefined\else\footerfont\@journal\hspace{8pt}\fi}
    }
}
\makeatother

\begin{document}

    \maketitle
    \thispagestyle{firststyle}

%----------------------------------------------------------

\noindent\textbf{Submission Overview.} \pending{confirm on Canvas how source code is submitted alongside this PDF.} The source package holds the C++17 code (shared core, serial, MPI+OpenMP and GPU drivers), the validation and measurement scripts, the experiment plans, the PBS job scripts used on the Vanda cluster, and a \texttt{README.md} with build and run instructions. As required, it contains no log files or terminal output; run records are kept separately to check the results in this report.

\section{Introduction}
\rhostart{T}he gravitational $N$-body problem asks for the motion of $N$ point masses under their mutual attraction. Direct summation costs $O(N^2)$ per force evaluation, which limits it to small $N$. The Barnes--Hut algorithm~\cite{BarnesHut1986} reduces the cost to $O(N \log N)$ by grouping distant particles into the cells of an octree and replacing each sufficiently distant group by its total mass at its centre of mass.

Barnes--Hut is a useful parallelization exercise because it is not embarrassingly parallel. Every particle's force depends on a tree built from all positions, so each timestep needs global information; the work per particle depends on local density, so equal particle counts need not mean equal work; and particles move, so any decomposition changes over time. We implement a serial reference and an MPI+OpenMP version in C++17, verify that the parallel version reproduces the serial one, and study how the split between MPI ranks and OpenMP threads, and the OpenMP loop schedule, affect scaling and efficiency on the Vanda cluster.

CPU performance has been measured only on 16 cores of one shared Vanda node, in one job with three repetitions. There, the force traversal takes at least 88\% of every strong- and weak-scaling run, and the position exchange at most 3\%. One rank with 16 threads ran 1.6 times faster than 16 single-thread ranks, and repetitions of the same case differed by up to 61\%. \textbf{No result at 64 or 256 cores exists yet}, so the required 256-core scaling is still missing. \pending{update once the multi-node jobs have run.}

\section{Problem and Method}
\subsection{Terms Used}
\begin{description}[noitemsep,leftmargin=0pt,font=\normalfont\itshape]
\item[Octree:] a tree of cubes in which each cell splits into eight equal octants.
\item[Leaf size $L$:] the most particles a cell may hold before it is split.
\item[Opening angle $\theta$:] the accuracy setting of Barnes--Hut. A larger value treats more cells as single masses, which is faster and less accurate.
\item[Monopole:] a cell's total mass placed at its centre of mass.
\item[Softening $\varepsilon$:] a small length added to each separation so that the force stays finite.
\item[Rank, thread, core:] an MPI process, an OpenMP thread within it, and a physical CPU core. $R \times T$ means $R$ ranks with $T$ threads each.
\item[Allocated cores $P$:] the cores a run holds, normally $R \times T$.
\item[Replicated tree:] every rank stores all positions and builds the same full tree.
\item[Owned particles:] the particles a rank computes forces for and moves.
\item[Binding, placement:] restricting a rank or thread to given cores, and the resulting layout.
\item[Strong and weak scaling:] more cores with fixed $N$, and more cores with $N/P$ fixed.
\item[Timed region:] the steps whose wall time is reported; setup and the warm-up steps before them are excluded.
\item[Traversal imbalance:] the slowest rank's force time divided by the mean over ranks.
\item[Manifest:] a file listing the runs of an experiment, fixed before they run.
\item[Gate:] a checked step in a job, with its own time limit.
\item[Schedule, chunk:] how OpenMP shares loop iterations among threads, and how many it hands out at a time.\item[NUMA domain:] a group of cores sharing one memory controller.
\item[Offload, device:] running the force loop on a GPU, which is the device.
\item[Kernel, flattening:] the force loop on its own, and copying the tree into plain arrays the device can read.
\item[PBS, \texttt{tm}, \texttt{pbs\_tmrsh}:] the cluster's job scheduler, its interface for starting processes, and its tool for starting a command on a node inside a job.
\item[Daemon:] the helper process Open MPI starts on a node to launch ranks.
\end{description}

\subsection{Physical Model and Inputs}
Particle $i$ has mass $m_i$, position $\mathbf{x}_i$ and velocity $\mathbf{v}_i$. With $G=1$ and Plummer softening $\varepsilon = 0.025$, its acceleration is
\begin{equation}
\mathbf{a}_i = \sum_{j \ne i} m_j \frac{\mathbf{x}_j - \mathbf{x}_i}{\left(|\mathbf{x}_j - \mathbf{x}_i|^2 + \varepsilon^2\right)^{3/2}}.
\label{eq:force}
\end{equation}
Softening keeps the force finite when particles come close or coincide.

Inputs are synthetic and generated by one routine, so every implementation starts from the same state. A fixed-seed Mersenne Twister supplies all random numbers. For each particle in ID order:
\begin{itemize}[noitemsep]
\item \emph{Mass:} every particle has mass $1/N$, so total mass is 1.
\item \emph{Radius:} a uniform random $m \in (0,1)$ gives $r = (m^{-2/3}-1)^{-1/2}$. This inverts the cumulative mass function of the Plummer density $\rho(r) \propto (1+r^2)^{-5/2}$ with scale length 1~\cite{Plummer1911,Aarseth1974}. The radius is not truncated, so rare particles lie far out and enlarge the root cell.
\item \emph{Direction:} uniform on the sphere, from a uniform $\cos\vartheta \in [-1,1]$ and azimuth $\in [0,2\pi)$.
\item \emph{Velocity:} speed uniform in $[0, 0.1]$, in an independent uniform direction.
\end{itemize}
These velocities are \emph{not} the Plummer equilibrium distribution, and we do not claim the cloud is in equilibrium: it starts cold and will evolve. Mean momentum is not removed. Our runs are a few steps long and measure cost, not astrophysical behaviour, so they do not need equilibrium. The profile is centrally concentrated, which makes tree depth and per-particle work vary across space; a uniform cube would hide the load-balance question. Synthetic data suit weak scaling because $N$ can be set freely with the same distribution. The cloud keeps the same size as $N$ grows, so density and tree depth grow too.

Time integration uses kick--drift--kick leapfrog~\cite{Verlet1967}: half kick $\mathbf{v} \leftarrow \mathbf{v} + \tfrac{\Delta t}{2}\mathbf{a}$, drift $\mathbf{x} \leftarrow \mathbf{x} + \Delta t\,\mathbf{v}$, new forces, half kick. Unless stated otherwise $\Delta t = 10^{-3}$.

\subsection{Barnes--Hut Force Evaluation}
The root cell is the smallest cube containing all particles, enlarged by a relative $10^{-9}$ so that boundary particles lie strictly inside. A cell holding more than $L$ particles (the \emph{leaf size}) is split into eight octants; particles are partitioned by a counting sort over the cell's index range, so every subtree owns a contiguous range of a permutation array. Subdivision stops at depth 64. Particles that still share a cell there, which happens only for coincident or near-coincident positions, remain together in one leaf, so no particle is lost. Each cell stores its total mass and centre of mass.

To evaluate $\mathbf{a}_i$ the tree is walked from the root. A cell of edge length $s$ whose centre of mass is at distance $d$ from $\mathbf{x}_i$ is accepted as a single softened point mass when $s/d < \theta$. Otherwise it is opened: an internal cell's children are visited, and an opened leaf is summed particle by particle with Eq.~\eqref{eq:force}, skipping $j=i$. One further rule is enforced: a cell that contains particle $i$ is never accepted, whatever $s/d$ is. For $\theta$ above about $1/\sqrt{3}$ such a cell can pass the test, and its monopole would include $i$'s own mass. This rule removes only that failure. Salmon and Warren~\cite{SalmonWarren1994} show that the standard criterion can give large worst-case errors for other configurations, which we do not address. At $\theta = 0$ no cell is accepted and the method reduces to direct summation, which provides an exact correctness check.

\subsection{Serial Baseline}
The serial program builds the tree and then evaluates every particle's acceleration in a single loop. Its timings are the denominator for every speedup, so it must not be artificially slow. Bucketed leaves reduce the number of tree nodes compared with one particle per leaf (867 against 2979 at $N=2000$). % Generated by code/scripts/report_tables.py from leafsweep-20261028T103743Z
In a local serial sweep ($N=50000$, 2 timed steps, median [min, max] of 3 runs, Docker container) the timed region took $L=1$: 3.68\,s [3.66, 3.71]; $L=4$: 3.30\,s [3.28, 3.55]; $L=8$: 3.10\,s [3.10, 3.28]; $L=16$: 3.01\,s [2.99, 3.66]; $L=32$: 3.29\,s [3.24, 3.41]. The lowest median, $L=16$, was 18\,\% below single-particle leaves, and its range overlapped that of $L=8$ and $L=32$.
 We use $L=8$ by default; \pending{the leaf size will be re-tuned on the cluster before the experiment manifest is frozen}. Direct summation is kept only as a correctness reference on small inputs and is never used as a speedup denominator.

\section{Parallel Implementation}
\subsection{Design: Replicated Tree, Distributed Targets}
We use the simplest decomposition that keeps the tree exact: every rank holds all positions and masses and builds the full tree, but each rank computes forces for, and integrates, only the particles it owns. Ownership is by contiguous particle-ID range. The first $N \bmod R$ of $R$ ranks own $\lceil N/R \rceil$ particles and the rest $\lfloor N/R \rfloor$; ranks with no particles are allowed and tested. Algorithm~\ref{alg:step} shows one timestep.

\begin{algorithm}[t]
\caption{One timestep on rank $r$ (owned IDs $[b_r, e_r)$)}
\label{alg:step}
\begin{algorithmic}[1]
\State $\mathbf{v}_i \gets \mathbf{v}_i + \tfrac{\Delta t}{2}\mathbf{a}_i$, \ $\mathbf{x}_i \gets \mathbf{x}_i + \Delta t\,\mathbf{v}_i$ \quad for $i \in [b_r, e_r)$
\State \texttt{MPI\_Allgatherv}: owned $\mathbf{x}$ $\to$ global $\mathbf{x}$ in ID order
\State build the global octree from all $N$ positions
\State \textbf{parallel for} $i \in [b_r, e_r)$ \texttt{schedule(runtime)}: $\mathbf{a}_i \gets$ tree walk
\State $\mathbf{v}_i \gets \mathbf{v}_i + \tfrac{\Delta t}{2}\mathbf{a}_i$ \quad for $i \in [b_r, e_r)$
\end{algorithmic}
\end{algorithm}

Rank 0 generates the initial state and broadcasts positions, velocities and masses once. Masses never change and are not resent. Velocities and accelerations stay distributed during the run and are gathered only for validation output. Particle coordinates are sent as three \texttt{MPI\_DOUBLE}s per particle; the code asserts at compile time that its vector type has exactly that layout. Counts and displacements are checked to fit in an \texttt{int}.

This design is chosen as a correct communication baseline, not as the most scalable option. Its costs are known in advance: every rank stores $O(N)$ positions, so memory per rank does not shrink with more ranks; every rank repeats the whole tree build; and the all-to-all position exchange moves $O(N)$ data per rank per step. A distributed, locally essential tree in the style of Warren and Salmon~\cite{WarrenSalmon1993} would avoid the replication, but it is outside the scope of this report.

\subsection{OpenMP Work Sharing and Memory Safety}
Within a rank, the tree is built by the main thread and is read-only during the force loop. Each iteration writes only $\mathbf{a}_i$ for its own $i$, keeps its accumulator and traversal stack in local variables, and reads shared data only. The loop therefore needs no locks or atomics. The traversal stack is a fixed local array whose size is proven sufficient: each visited cell pushes at most eight children and depth is capped at 64, so at most $7 \cdot 64 + 8$ entries are ever pending.

Traversal cost varies strongly between particles, since particles in the dense core open more cells. The loop uses \texttt{schedule(runtime)}, so static, dynamic and guided schedules and their chunk sizes can be compared on the same binary~\cite{OpenMP51}. Both the requested schedule and the schedule the runtime reports are recorded with every run.

\subsection{MPI and Threads}
MPI is initialised with \texttt{MPI\_Init\_thread} requesting \texttt{MPI\_THREAD\_FUNNELED}, and the program stops if the library provides less~\cite{MPI40}. All MPI calls are made by the main thread outside parallel regions. Argument errors are detected identically on every rank, so all ranks exit together rather than leaving some waiting in a collective; an unexpected exception on one rank aborts the whole job. After the timed loop, every rank checks that its global positions and masses are bitwise identical to every other rank's and that owned counts sum to $N$; a mismatch makes the run fail.

Placement matters as much as the code. During local development we found that Open MPI's default binding can put a rank's whole thread team on one hardware thread: % Generated by code/scripts/report_tables.py from binding-20261028T103451Z
With one rank and 4 OpenMP threads ($N=20000$, 3 timed steps, 3 repetitions each, local Docker container), Open MPI's default placement gave a median timed region of 1.57\,s against 0.40\,s with \texttt{--bind-to none}. Open MPI's \texttt{--report-bindings} showed the default binding the rank, and therefore its whole thread team, to a single hardware thread.
 Every measurement run and placement probe therefore sets placement explicitly: $R$ ranks per node with $T$ cores each (Open MPI \texttt{--map-by ppr:$R$:node:PE=$T$ --bind-to core}), and within each rank one OpenMP thread per core (\texttt{OMP\_PROC\_BIND=close}, \texttt{OMP\_PLACES=cores}). A placement probe run inside each job checks that every thread runs on a different CPU and is bound to a single core. Placement is recorded with every run and treated as part of the configuration in the analysis.

\section{Correctness}
Each test below exists to catch a particular fault; Table~\ref{tab:test-purpose} gives the reason for each. We separate two questions: how far Barnes--Hut is from the exact force at a given $\theta$ (approximation error), and whether the parallel program computes the same Barnes--Hut result as the serial one (implementation error). The approximation and implementation results below come from a correctness job on a Vanda compute node (PBS job 1503779, 8 allocated cores, compiled there with GCC~13.2 and Open MPI~4.1.6, with hybrid--serial comparisons run unbound); the same gates also passed earlier in a local Debian container with GCC~12 and Open MPI~4.1.4. The energy study (Table~\ref{tab:energy}) and the mutation check are local development evidence only.

\begin{table}[t]
\centering
\small
\caption{Why each validation test is meaningful.}
\label{tab:test-purpose}
\begin{tabular}{@{}p{0.40\columnwidth}p{0.54\columnwidth}@{}}
\toprule
Test & Fault it would reveal \\
\midrule
$\theta=0$ against direct summation & A particle lost or counted twice in the tree, or self-interaction \\
Tree invariants & A particle in no leaf or two; wrong cell mass \\
Two particles, closed form, $\theta$ up to 10 & Wrong force law; a cell accepted that contains its target \\
Coincident and extreme inputs & Particles dropped at the depth limit \\
Hybrid against serial, by particle ID & Missing exchange, wrong ownership, a race \\
Uneven $N$, empty ranks, $N=1$ & Errors in dividing particles between ranks \\
Repeated threaded runs & Results that depend on thread timing \\
Perturbed force, mass mismatch, removed exchange & A test harness that cannot fail \\
Invalid arguments & Ranks left hanging in a collective \\
Placement probe & Threads sharing a core, which would spoil timings \\
GPU against direct summation and serial & A missing data transfer; silent fallback to the CPU \\
\bottomrule
\end{tabular}
\end{table}

Correctness across two nodes rests on one job; every earlier two-node job stopped before the comparison. Job 1506687 stopped at its pre-build allocation check, which at that time required identical CPU models and was given two different ones. Job 1507215 confirmed two distinct hosts, the 8-core allocations and data exchange between nodes, then stalled in a placement probe and reached its time limit. A diagnostic job (1514029) reproduced the stall: under PBS's own process manager, launches that used a core mapping and binding hung, even for \texttt{hostname}, while one launch through PBS's \texttt{pbs\_tmrsh} agent passed. We did not find the cause. Multi-node jobs therefore start their remote daemons through \texttt{pbs\_tmrsh}, with no fallback to ssh, and check that those daemons run under the PBS daemon on every node. % Generated by code/scripts/report_tables.py from twonode-20261103T102750Z-1514458
PBS job 1514458 (commit \texttt{0b1f36b}; nodes IN-002 and IN-005, both Intel Xeon Gold 6438N; 8 allocated cores each) passed 17 of 17 critical and required steps. In its hybrid--serial step, 22 of 22 cases were within tolerance, and the largest difference was 0. 17 of the cases ran across both nodes, with up to 16 ranks and $N=20011$, all with binding off (\texttt{--bind-to none}). 8 of 8 bound placement probes passed.
 Single-node measurement jobs at first used a plain node file instead. Pilot job 1519637 showed that this did not avoid \texttt{tm}: \texttt{mpirun} still selected it, started a daemon on its only node and made no progress at the first bound launch, with no run started and nothing measured. The process state matched the earlier stall, so we take it to be the same fault, though that is an inference. All measurement jobs now launch through \texttt{pbs\_tmrsh}. \pending{whether this works on one node at full size is not yet shown on Vanda.}

These results answer two separate questions. The comparison against the serial program shows that the distributed computation gives the same answer across two nodes, but it ran unbound, so it says nothing about bound runs. The bound runs are the placement probes, which check where threads run and do not run the solver, and one short solver run with 4 ranks of 4 threads ($N=50000$, 2 steps), which checks that a bound two-node run completes and is recorded, without comparing its result with serial. No bound result has been compared with serial. This is also one job on one pair of nodes, at far below full scale: it does not show that the launch method is reliable, explain the stall, or cover the larger allocations and other node types used for measurement.

\subsection{Approximation Error}
Table~\ref{tab:approx} compares serial Barnes--Hut with direct summation on the same input. At $\theta=0$ agreement is at rounding level, which tests tree construction and traversal: every particle is reached exactly once and self-interaction is excluded. At $\theta>0$ the errors measure the approximation and grow with $\theta$; they are diagnostics, not pass/fail thresholds. The acceptance threshold at $\theta=0$ is a maximum relative error of $10^{-10}$.

% Generated by code/scripts/report_tables.py from artifacts/cluster/runs/correctness-20261028T123124Z-1503779/serial_validation_vanda-1503779.json
\begin{table}[t]
\centering\small
\caption{Approximation error of serial Barnes--Hut against direct summation, $N=2000$, seed 42, relative error of the acceleration vector per particle. Vanda compute node, PBS job 1503779 (Intel Xeon Gold 6438N, Open MPI~4.1.6, GCC~13.2), 8 allocated cores.}
\label{tab:approx}
\begin{tabular}{rrrr}
\hline
$\theta$ & leaf & max rel. & mean rel. \\
\hline
0.0 & 1 & $1.1\times10^{-14}$ & $1.5\times10^{-15}$ \\
0.3 & 1 & $1.6\times10^{-2}$ & $8.6\times10^{-4}$ \\
0.5 & 1 & $1.1\times10^{-1}$ & $3.5\times10^{-3}$ \\
0.0 & 8 & $1.1\times10^{-14}$ & $1.5\times10^{-15}$ \\
0.3 & 8 & $8.4\times10^{-3}$ & $6.7\times10^{-4}$ \\
0.5 & 8 & $5.9\times10^{-2}$ & $3.0\times10^{-3}$ \\
\hline
\end{tabular}
\end{table}


Edge cases are tested by a built-in self-test that checks tree invariants (each particle in exactly one leaf, root mass equal to total mass) and $\theta=0$ agreement with direct summation for each case at leaf sizes 1 and 8, plus closed-form and symmetry checks. % Generated by code/scripts/report_tables.py
18 of 18 self-test checks passed, covering: all\_coincident, below\_max\_depth, coincident\_cluster, empty, lattice\_on\_planes, negative\_control\_force\_perturbation, negative\_control\_mass\_mismatch, singleton, symmetric\_cube, two\_particle\_analytic, unequal\_masses.
 The two-particle case uses the closed-form softened force at $\theta=0$, $0.5$ and $10$. At $\theta=10$ the root passes $s/d<\theta$ from either particle, so this case fails if a cell containing the target is ever accepted; we confirmed that by temporarily restoring that behaviour. The negative controls confirm that the checks can fail: a $10^{-8}$ perturbation of one force and a mass mismatch must both be detected. The serial code also runs cleanly under AddressSanitizer and UndefinedBehaviorSanitizer.

\subsection{Implementation Error}
The hybrid program is compared with the serial one at identical $\theta$, leaf size, softening, timestep, seed and particle count, by particle ID. Tolerances were fixed before any comparison was run (Table~\ref{tab:hybrid}) and are not adjusted afterwards. The cases cover 1--5 ranks, 1--4 threads, all three schedules, $N$ not divisible by the number of ranks, ranks that own no particles, and $N=1$.

% Generated by code/scripts/report_tables.py from hybrid_validation_linux.json
\begin{table}[t]
\centering\small
\caption{Implementation error: MPI+OpenMP against serial Barnes--Hut at identical settings ($\theta=0.5$ unless the case tests otherwise), largest absolute difference over positions, velocities and accelerations by particle ID. Tolerances: one step $1e-12|a|+1e-14\max|a|$; ten steps $1e-10|x|+1e-13\max|x|$. Local development run (Docker, Open MPI~4.1.4, GCC~12), not cluster evidence.}
\label{tab:hybrid}
\begin{tabular}{llrrrl}
\hline
ranks$\times$thr. & schedule & $N$ & steps & max diff & \\
\hline
1$\times$1 & static/0 & 2000 & 1 & 0 & pass \\
1$\times$4 & static/0 & 2000 & 1 & 0 & pass \\
1$\times$4 & dynamic/1 & 2000 & 1 & 0 & pass \\
2$\times$1 & static/0 & 2000 & 1 & 0 & pass \\
3$\times$2 & dynamic/16 & 2000 & 1 & 0 & pass \\
4$\times$2 & guided/0 & 2000 & 1 & 0 & pass \\
3$\times$2 & static/0 & 1500 & 1 & 0 & pass \\
2$\times$2 & static/0 & 1500 & 1 & 0 & pass \\
1$\times$1 & static/0 & 2000 & 10 & 0 & pass \\
1$\times$4 & dynamic/8 & 2000 & 10 & 0 & pass \\
3$\times$2 & static/0 & 2001 & 10 & 0 & pass \\
4$\times$2 & guided/0 & 2003 & 10 & 0 & pass \\
4$\times$1 & static/0 & 7 & 10 & 0 & pass \\
5$\times$1 & static/0 & 3 & 10 & 0 & pass \\
3$\times$1 & static/0 & 1 & 5 & 0 & pass \\
\hline
\end{tabular}
\end{table}


Every difference in Table~\ref{tab:hybrid} is exactly zero. That is expected rather than suspicious: each particle's force is computed by the same code on a bitwise-identical tree in the same order, whatever rank or thread computes it. To check that the harness can detect a real fault, we removed the position exchange from a test build. All ten multi-rank cases then failed, as did the repeatability check, while single-rank cases still passed, as they should. % Generated by code/scripts/report_tables.py
Three repeated runs with 2 ranks $\times$ 4 threads and dynamic chunk 1 were bitwise identical; a $10^{-9}$ relative perturbation of one position was detected by the comparator; evidence recorded 2026-10-28; sources \texttt{nbody.hpp} 7f8fe48d, \texttt{run\_support.hpp} cfedf58d, \texttt{bh\_hybrid.cpp} ea8df96b, \texttt{barnes\_hut.cpp} 562b0fa8.
 Invalid arguments make all ranks exit with status 2 without hanging.

\subsection{Energy Behaviour}
Energy is a supplementary diagnostic, not a correctness proof (Table~\ref{tab:energy}). With exact forces ($\theta=0$) the relative energy change falls by a factor of four each time $\Delta t$ is halved, the second-order behaviour expected of leapfrog. For $\theta>0$ it is almost independent of $\Delta t$ and instead differs between opening angles: it increases from $\theta=0.2$ to $0.5$ but is slightly lower at $0.7$ than at $0.5$. A plateau that depends on $\theta$ but not on $\Delta t$ is consistent with a contribution from the approximate tree force. Two further facts complicate the picture: the energy is evaluated exactly while the dynamics use approximate forces, and the tree changes discontinuously between steps. We have not isolated the cause, and draw no conclusion beyond the observed values.

% Generated by code/scripts/report_tables.py from energy_study_linux.json
\begin{table*}[t]
\centering\small
\caption{Relative change in total energy over simulated time $T=0.2$ ($N=2000$, leaf size 8). Energy is evaluated exactly ($O(N^2)$), outside any timed region. Local development run, supplementary diagnostic only.}
\label{tab:energy}
\begin{tabular}{rrrrrr}
\hline
$\theta$ & $\Delta t=0.008$ & $\Delta t=0.004$ & $\Delta t=0.002$ & $\Delta t=0.001$ & $\Delta t=0.0005$ \\
\hline
0.0 & $7.3\times10^{-8}$ & $1.8\times10^{-8}$ & $4.5\times10^{-9}$ & $1.1\times10^{-9}$ & $2.8\times10^{-10}$ \\
0.2 & $1.2\times10^{-7}$ & $1.8\times10^{-7}$ & $1.9\times10^{-7}$ & $1.9\times10^{-7}$ & $1.9\times10^{-7}$ \\
0.3 & $1.0\times10^{-6}$ & $1.1\times10^{-6}$ & $1.1\times10^{-6}$ & $1.1\times10^{-6}$ & $1.1\times10^{-6}$ \\
0.5 & $3.6\times10^{-6}$ & $3.7\times10^{-6}$ & $3.7\times10^{-6}$ & $3.7\times10^{-6}$ & $3.7\times10^{-6}$ \\
0.7 & $3.5\times10^{-6}$ & $3.6\times10^{-6}$ & $3.6\times10^{-6}$ & $3.6\times10^{-6}$ & $3.6\times10^{-6}$ \\
\hline
\end{tabular}
\end{table*}


\section{Computing Environment}
Formal measurements use the NUS HPC cluster Vanda, which is scheduled by PBS. Code is compiled on the allocated compute node, because login and compute nodes differ in architecture, with \texttt{-O3 -march=native} and the same flags for serial and hybrid builds. Jobs are submitted to the zero-rate \texttt{auto\_free} routing queue, which places multi-core jobs in \texttt{cpu\_parallel}. The first correctness job ran on node IN-003: two Intel Xeon Gold 6438N sockets with 32 cores each, one hardware thread per core, two NUMA domains and Red Hat Enterprise Linux~9.4. It used the \texttt{OpenMPI/4.1.6-GCC-13.2.0} module (GCC~13.2.0, Open MPI~4.1.6). PBS confined the job to its 8 allocated cores, and the placement probe confirmed one thread per core.

The free node pool is not uniform: besides the 64-core Gold~6438N (IN) nodes it contains 72-core Xeon Platinum~8452Y (CN) nodes, and the scheduler offers no resource that selects a CPU model, so a multi-node request can receive both types (job 1506687 did). Each run records its nodes and CPU models. Timings are compared only within a single CPU model, and a speedup's serial baseline must come from the same model. Multi-node correctness jobs, which may receive mixed nodes, use one portable build (\texttt{-march=x86-64 -mtune=generic}, recorded together with the target the compiler resolved) instead of \texttt{-march=native}; their results are correctness evidence only. The free queue allows at most 160 CPUs per job and per user, so 256 cores needs the charged \texttt{batch\_cpu} queue. \pending{whether a 256-core run is approved and completed: not yet decided, and the report must state any shortfall.} \pending{memory per node, interconnect, how measurement jobs are restricted to one node type, and the nodes and allocations used by the measurement jobs.}

Local development used a Debian~12 container (GCC~12.2, Open MPI~4.1.4) on an 8-core Intel Core i9-9880H laptop with simultaneous multithreading, and Apple clang on macOS for the serial gates. No timing from these machines is reported as a performance result.

\section{Experimental Methodology}
\subsection{Timing}
The timed region is a fixed number of timesteps after setup. Setup covers initial-condition generation (and its broadcast in the hybrid code), the initial force evaluation, optional diagnostics and any warm-up steps, and is recorded separately. In the hybrid code all ranks synchronise with a barrier, time the region with \texttt{MPI\_Wtime}, and the reported time is the maximum over ranks. Within the region each rank accumulates four phases (integration, position exchange, tree build and force traversal) and each is reported as its maximum over ranks. Because different ranks can be slowest in different phases, these maxima need not sum to the wall time. The traversal imbalance is the ratio of the maximum to the mean traversal time over ranks. The serial program times the same region and phases with a monotonic clock. Energy diagnostics and output are never inside the timed region.

\subsection{Metrics}
$P$ is the number of physical cores allocated to a run, normally ranks $\times$ threads, and runs that place more threads than allocated cores are rejected. Strong-scaling speedup is $S(P) = T_{\mathrm{serial}}/T(P)$ and efficiency is $E(P) = S(P)/P$, where $T_{\mathrm{serial}}$ is serial Barnes--Hut with the same problem, numerical settings, source revision, compiler flags and CPU model. The one-rank, one-thread hybrid time is reported alongside it to expose MPI and OpenMP overhead. Weak-scaling efficiency is $T(P_0, N_0)/T(P, N)$ with $N/P = N_0/P_0$ and a stated baseline $P_0$. Constant particles per core does not mean constant work per core, since tree depth grows with $N$, and the replicated design stores all $N$ positions on every rank.

\subsection{Planned Studies}
None of these CPU studies has run yet. They are numbered as in the experiment plans.
\begin{enumerate}[noitemsep]
\setcounter{enumi}{-1}
\item Serial runtime against particle count $N$.
\item Strong scaling at fixed $N$ and step count, to at least 256 cores, including serial and one-rank, one-thread baselines, with speedup and efficiency.
\item Ranks $\times$ threads at fixed $P$ and $N$, e.g.\ $16{\times}1$, $8{\times}2$, $4{\times}4$, $2{\times}8$, $1{\times}16$, with matched placement.
\item OpenMP schedule: static against dynamic and guided, with a small chunk-size sweep fixed in advance.
\item Weak scaling with $N \propto P$ and fixed threads per rank, to at least 256 cores.
\item Leaf size, on the serial program.
\item Accuracy against cost across $\theta$, reported separately from scaling.
\end{enumerate}
The share of time in the position exchange is read from the strong and weak runs. Axes for time, particle count and CPU count are logarithmic in every figure; efficiency and the exchange share are ratios and use linear vertical axes.

Problem sizes, step counts and core counts will be chosen from pilot runs within the cluster's memory and time limits and then frozen in a manifest: a machine-readable file in the code repository whose identifier and checksum every run record carries. Formal measurement jobs refuse a manifest that is not frozen. Two one-repetition pilots, on 16 cores and on a whole 64-core node, set the sizes: $N=10^6$ for strong scaling, 50{,}000 particles per core for weak scaling, $\theta=0.5$, leaf size 8, dynamic schedule with chunk 64, 3 timed steps after 1 warm-up step. They are not reported as results. The 16-core study is frozen as \texttt{studies-16core-v1}. \pending{frozen manifests for 64 to 256 cores.} Each configuration is run at least three times. We report the median with the minimum and maximum, and failed runs are kept and listed.

\subsection{Bounded Runs and Failed Runs}
A step that never finishes must not consume a whole allocation silently, as the stalled placement probe did in job 1507215. Every step of a correctness or diagnostic job now has its own time limit, about ten times its observed duration. Before each step the limit is shortened, if necessary, so that the step's worst case, including the time to stop it, record process states and check every node, still ends before the allocation does; a step with too little time left is skipped. This is a bounded best effort rather than a guarantee: work outside the steps is not limited, and a process that cannot be killed would still run on. Shortly before a limit expires, the state of every process on every node is recorded. After every step, processes it left behind are terminated on every node, but only those that can be shown to belong to the job, and their absence is checked. If that check fails, or an MPI process of unknown ownership remains, the job stops, because later steps might share the cores with stray processes. Steps that depend on a failed step are skipped. Each step is classed as critical, required or optional; a run passes only if every critical and required step passes, while optional comparisons are listed without deciding the outcome. Timed-out, failed and skipped steps are kept in the record. Measurement jobs run the same way, and within them every timed run has its own time limit; a run that could not finish before the job's deadline is recorded as skipped rather than started.

\subsection{Provenance}
Each run writes its own record: exact command, scheduler job ID, nodes, CPU model, allocated cores, OpenMP settings, source revision with a hash of the source tree, binary hash, compiler, flags and MPI library, raw output and the program's result. Summaries, tables and figures are regenerated from these records by scripts, never copied by hand. Records are grouped by every setting that can change a timing, including placement and build, before medians are taken.

\section{Results}
% Each fragment below is written by code/scripts/report_figures.py from
% analysed Vanda runs (artifacts/cluster/) only; until it exists the slot
% shows a pending marker. Do not type numbers into these slots by hand.
\newcommand{\resultslot}[2]{\IfFileExists{generated/#1.tex}{\input{generated/#1.tex}}{\pending{#2}}}
Every CPU result below comes from one job, PBS 1523703 (commit \texttt{e86bf5d}), on 16 cores of node IN-002, an Intel Xeon Gold~6438N. All 9 of its correctness and placement steps passed, including a bound comparison with the serial program on these cores, and all 121 runs completed (Table~\ref{tab:failed-runs}). Each case ran three times, with repetitions interleaved: the first of every case, then the second, then the third. The other 48 cores of the node were not ours, so other jobs may have shared memory bandwidth and cache with ours; the records cannot show this. All 16 cores were on one socket. \pending{no result yet at 64 or 256 cores: the required 256-core strong and weak scaling are missing.} Statements marked ``expected'' follow from the design and are not results.

\paragraph{How much the runs varied.} For a given case the slowest repetition took up to 61\% longer than the fastest, and 18\% longer in the median case. The third repetition was slower than the first in 27 of 39 cases. Serial and hybrid $1\times1$ run the same algorithm on one core, yet their ratio ranged from 0.95 to 1.30 across repetitions. Speedup against serial therefore depends on which baseline run it is divided by. Table~\ref{tab:repetitions} shows each repetition against both baselines, and the figures show the range as bars.

\subsection{Serial Runtime against Problem Size}
\resultslot{fig_serial_n}{serial runtime against $N$, log--log axes (study 0).}
Figure~\ref{fig:serial-n} spans $N=10^4$ to $2\times10^6$. Before measuring we expected about $O(N\log N)$ per step, which over this range would give a log--log slope of about 1.09. The measured slope is 1.33, so time per particle grew faster than $\log N$: traversal took about 9.5\,$\mu$s per particle per step at $N=10^4$ and about 50\,$\mu$s at $2\times10^6$. We did not test why. Plausible contributors, untested, are a working set that outgrows the caches and the shared node's conditions. Repetitions agree within 1\% up to $N=3\times10^4$, and differ by up to 18\% from $N=10^5$ upwards.

\subsection{Strong Scaling}
\resultslot{fig_strong_runtime}{strong-scaling runtime figure to at least 256 cores (study 1).}
\resultslot{fig_strong_speedup}{strong-scaling speedup figure (study 1).}
\resultslot{fig_strong_efficiency}{strong-scaling efficiency figure (study 1).}
At $N=10^6$ and 16 cores, MPI-only efficiency against the serial median was 0.87, and one rank with 16 threads reached 1.42 (Figure~\ref{fig:strong-efficiency}). Against hybrid $1\times1$ of the same repetition, the ranges are 0.74--0.95 for $16\times1$ and 1.12--1.21 for $1\times16$ (Table~\ref{tab:repetitions}).

\emph{Efficiency above 1.} Efficiency above 1, or superlinear speedup, means the parallel run used less total core-time than the baseline. This is possible and has known causes, for example a working set that fits in the combined caches of many cores but not in one core's cache. Here part of it comes from the baseline: the serial runs were slower than hybrid $1\times1$ in two of three repetitions. It does not all come from the baseline, though, because against $1\times1$ the thread-only runs stay above 1 in every repetition ($1\times4$: 1.06--1.27, $1\times8$: 1.14--1.26, $1\times16$: 1.12--1.21). We did not measure cache behaviour, so we do not know the cause. A one-repetition sizing run on a whole 64-core node, which is not reported as a result, showed no efficiency above 1. So this should not be read as a property of the program at larger core counts.

Expected before measuring: the tree build and the exchange do not shrink with more ranks, so efficiency should fall as $P$ grows. At 16 cores neither explains much: the build took at most 10\% of wall time and the exchange at most 1.2\% in the strong runs (Figures~\ref{fig:phases} and~\ref{fig:comm-fraction}). \pending{strong scaling from 16 to at least 256 cores.}

\subsection{Ranks and Threads}
\resultslot{tab_ranks_threads}{ranks $\times$ threads table (study 2).}
At 16 cores, fewer ranks with more threads were faster: the median went from 5.6\,s at $1\times16$ to 9.2\,s at $16\times1$ (Table~\ref{tab:ranks-threads}). That order holds at the extremes, but $8\times2$ was slower than $16\times1$, so it is not monotonic. Most of the gap is in the traversal itself (8.5\,s against 5.1\,s), not in the exchange (1.2\% of the $16\times1$ run) or the build. Imbalance between ranks was at most 1.2\%. One possible reason is that 16 separate tree copies compete for cache, where one shared copy serves 16 threads. We did not test it.

\subsection{OpenMP Schedule}
\resultslot{tab_schedule}{schedule and chunk table (study 3).}
Dynamic scheduling with chunk 64 had the lowest median at $1\times16$, by up to 22\%. Every other schedule's range overlaps its range, so three repetitions on a shared node do not separate them. At $4\times4$ static took a 10\% longer median than dynamic, again with overlapping ranges. We keep dynamic, chunk 64.

\subsection{Weak Scaling}
\resultslot{fig_weak_runtime}{weak-scaling runtime figure to at least 256 cores (study 4).}
\resultslot{fig_weak_efficiency}{weak-scaling efficiency figure (study 4).}
With 50{,}000 particles per core, MPI-only weak efficiency fell from 1 at one core to 0.28 at 16, and with 8 threads per rank it was 0.71 from 8 to 16 cores (Figure~\ref{fig:weak}). We expected a fall, since constant particles per core is not constant work: $N$ grows with $P$, so every rank's traversal deepens and its tree build covers all $N$. In these runs traversal still took over 88\% of time, the build rose from 0.7\% to 9.9\%, and the exchange stayed under 2\%. Memory per rank grows with $N$, not $N/P$, because positions and the tree are replicated: about 26\,MiB at $N=5\times10^4$ and 129\,MiB at $8\times10^5$ per rank. \pending{weak scaling from 16 to at least 256 cores.}

\resultslot{tab_repetitions}{per-repetition baseline table.}

\subsection{Where the Time Goes}
\resultslot{fig_phases}{phase times against core count.}
\resultslot{fig_comm_fraction}{share of runtime spent in the position exchange.}
Each step sends every particle's three coordinates to every rank, so each rank receives $24N$ bytes and a node with $R$ ranks receives $24NR$ bytes. At fixed $N$ the force time falls with $P$ while this does not, so we expected the exchange's share to rise with $P$. It did, but stayed small on one node: 0.2\% at 2 ranks and 1.2\% at 16 in the strong runs, and 1.9\% at 16 in the weak runs. All of this was within one node, through shared memory; between nodes the exchange crosses the network. \pending{communication between nodes, up to at least 256 cores.}

\subsection{GPU Test}
The GPU program uses the serial tree code on the CPU and offloads only the force traversal, with OpenMP \texttt{target}. A silent fallback to the CPU would make any GPU result false, so the job requires a device: the kernel must report that it ran on a device whose memory is separate from the host's, offload is set to mandatory, and \texttt{nvidia-smi} provides a view from outside the program. Separate memory also matters for correctness, because a missing data transfer cannot show up when the target region runs on the host. % Generated by code/scripts/report_tables.py from gpu_validate-20261103T134410Z-1514775
PBS job 1514775 (commit \texttt{fa9da68}, NVIDIA A40, \texttt{nvc++} 26.1-0) passed 23 of 23 checks. At $\theta=0$ the device force differed from direct summation by at most $5.0\times10^{-14}$ (relative), and a 20-step trajectory differed from the serial program by at most $9.7\times10^{-14}$. While a run of $N=10^6$ had offload set to mandatory, \texttt{nvidia-smi} reported up to 100\% GPU utilisation and listed the program as a process on the GPU.
 The deliberate-fault checks of the validator ran only in local development without a device, not on the GPU. This is one validation job on one GPU and shows correctness, not speed.

\resultslot{fig_gpu_runtime}{GPU timing against the CPU cores of the same allocation: not measured yet.}
\resultslot{tab_gpu}{device, CPU backend, kernel, flattening and transfer times.}
In this timing job (PBS 1519638, commit \texttt{8a50c8f}; Figure~\ref{fig:gpu-runtime} and Table~\ref{tab:gpu}), one A40 took longer than the 36 CPU cores of the same allocation at every size. The ratio of CPU to device time rose from 0.46 at $N=10^5$ to 0.74 at $N=3\times10^6$; below 1, the device is slower. The force kernel accounts for at least 97\% of the difference. Node mapping and both transfers together took under 1\% of device wall time, so moving data is not what limits this program. Tree build, flattening and integration use one host thread on both backends. We did not tune the kernel for the GPU and did not test why it is slower. We also ran nothing above $N=3\times10^6$, so the rising ratio does not show that the device would overtake the CPU at a larger size. The CPU backend is this same kernel built with \texttt{nvc++}, not the MPI+OpenMP program, so the comparison says nothing about that solver. These figures come from one job on one node with one GPU and three repetitions per case.

\subsection{Accuracy and Cost}
\resultslot{tab_accuracy_cost}{accuracy against cost table (study 6).}
\resultslot{tab_leaf}{leaf-size table (study 5).}
Lowering $\theta$ from 0.5 to 0.3 cut the maximum relative force error from $6.8\times10^{-2}$ to $1.4\times10^{-2}$, and took 4.3 times as long at $4\times4$. Raising it to 0.9 made the run 2.9 times faster, with a maximum error of 0.39. We keep $\theta=0.5$. Errors were measured at $N=20{,}000$ against direct summation, so they show the trend, not the error at $N=10^6$. For leaf size, 8 had the lowest median at $4\times4$, and 4, 16 and 32 took 16 to 33\% longer (Table~\ref{tab:leaf}). Serially at $3\times10^5$, 8 and 16 were equal and 32 was 11\% slower, within wide ranges.

\subsection{Failed and Skipped Runs}
No run of job 1523703 failed, timed out or was skipped (Table~\ref{tab:failed-runs}). Two earlier pilot jobs ended before any run and so measured nothing. Job 1514774 stopped at entry after one second because of a job-script path error. Job 1519637 stopped at its placement step (see Correctness). Two later pilots, jobs 1523275 (16 cores) and 1522998 (a whole 64-core node), completed all their runs. They were used only to choose sizes: each has one repetition, so their times are not reported as results. \resultslot{tab_failed_runs}{table of failed, timed-out and skipped runs.}

\section{Discussion}
\subsection{Parallelization Decisions}
Table~\ref{tab:decisions} lists the technique used for each part of the program and the communication it needs. These are design choices. Their cost has been measured on 16 cores only.

\begin{table}[!ht]
\centering
\small
\caption{Technique and communication for each part of a timestep.}
\label{tab:decisions}
\begin{tabular}{@{}p{0.2\columnwidth}p{0.38\columnwidth}p{0.3\columnwidth}@{}}
\toprule
Part & Technique & Communication \\
\midrule
Initial state & Rank 0 generates it & One broadcast \\
Integration & Each rank, owned particles only & None \\
Positions & Replicated on every rank & \texttt{MPI\_Allgatherv} each step \\
Tree build & Whole tree on each rank, one thread & None; the same tree results from the same positions \\
Force loop & OpenMP within a rank, ranks by ID range & None; the tree is read-only \\
GPU force loop & OpenMP \texttt{target} on one device & Host--device transfers, no MPI \\
\bottomrule
\end{tabular}
\end{table}

Replication makes the force loop free of communication and keeps results exactly equal to serial, which is what the correctness tests rely on. The price is the per-step all-to-all exchange and a tree build that is repeated on every rank. Threads avoid neither cost, but fewer ranks per node reduce the exchange volume and memory, which is why the ranks $\times$ threads study matters.

\subsection{What the Results Support}
The correctness jobs support the claim that the parallel program reproduces the serial result in the configurations tested, within the limits given in Correctness. The GPU job supports only the comparison in its section: one A40 was slower than 36 CPU cores with this kernel. Job 1523703 supports claims about 1 to 16 cores of one shared node: there, traversal dominates, the exchange costs under 3\%, and threads did better than ranks on the same cores. It does not support any claim about 64 or 256 cores, communication between nodes, or how the program behaves on a node it has to itself. Its speedups against serial depend on which baseline run is used, and Table~\ref{tab:repetitions} gives the range. \pending{claims at 64 to 256 cores, after the multi-node jobs.}

\subsection{Limitations}
\label{sec:limitations}
The replicated design keeps $O(N)$ memory and a full tree build on every rank, and exchanges $O(N)$ positions per step, so it cannot scale to very large $N$ or very many ranks. The tree build is serial within each rank. Ownership by particle ID ignores spatial locality, and neither locality-aware ordering nor work-weighted partitioning is implemented. Inputs are synthetic Plummer-profile clouds, so no conclusions about real astrophysical data are drawn. The measurements so far share one node with other jobs, and the same case varied by up to 61\%. Three repetitions show that spread, but they are too few to separate settings that differ by less. Sizing runs on other nodes, not reported as results, ran the same binaries faster per step, so absolute times here are specific to this allocation. \pending{limitations revealed by the multi-node measurements.}

\section{Conclusion}
The replicated MPI+OpenMP Barnes--Hut program reproduces the serial result, and on 16 cores of one node its time is dominated by the force traversal rather than by communication. Splitting those cores into threads rather than ranks was faster, for reasons we did not establish. Run-to-run variation on the shared node was large enough to change the speedup by tens of percent, so we report ranges rather than single values. \pending{conclusions on strong and weak scaling to at least 256 cores.}

\subsection{AI Usage Statement}
The group independently developed the experimental design and workflow, including the plans for data processing and results analysis. OpenAI Codex, GPT and Anthropic's Claude assisted with code implementation and review, environment configuration, data processing and analysis following the group's plans, experiment execution, result verification, and report drafting and revision.

%======================
% References
%======================

\printbibliography

\end{document}
